\sectioncentered*{Заключение}
\addcontentsline{toc}{section}{Заключение}

В рамках выполненной курсовой работы спроектирована, реализована и протестирована гибридная интеллектуальная система автоматической видеорегистрации и анализа шахматных партий, объединяющая нейробионический контур компьютерного зрения и логико-символьный решатель на базе формальных правил ФИДЕ. Все задачи, сформулированные во введении, выполнены в полном объёме.

По итогам работы получены следующие основные результаты:
\begin{enumerate}
    \item Проанализирована предметная область автоматической видеорегистрации шахматной игры, формализованы сущности и межмодульные контракты данных (\texttt{Frame}, \texttt{BoardGeometry}, \texttt{CellObservation}, \texttt{BoardObservation}, \texttt{MoveHypothesis}, \texttt{MoveEvent}, \texttt{PipelineState}), разделена ответственность между уровнями восприятия, временного дебаунсинга, символьной валидации, персистентности и аналитики.
    \item Спроектирован и реализован помехоустойчивый гибридный конвейер компьютерного зрения: двухстадийная локализация доски (\texttt{CoarseToFineLocalizer} с 1D гомографией поперечного отношения, робастной оценкой Хубера и трекером дрифта \texttt{BoardGeometryTracker}), не зависящее от яркости определение ориентации и полнокадровое распознавание позиции трансформером \texttt{ChessQ-Lite} (\texttt{ViT-S14 INT8}), дополненное семантическим детектором изменений \texttt{DINOv2-Small INT8}, структурным базисным ядром клетки (\textit{Base-Core}) и детектированием перекрытия рукой.
    \item Реализовано логико-символьное ядро с 9-состоянным конечным автоматом \texttt{GameFSM}, многофакторным скорером $S(m)$, шлюзом уверенности \texttt{DecisionGate}, ограниченным многоходовым восстановлением при быстрой игре, физическим откатом ходов (\textit{takeback}), строгой фильтрацией нелегальных перемещений фигур с фиксацией инцидентов \texttt{illegal\_attempt} и криптографическим журналом событий \texttt{events.jsonl} с SHA-256 связыванием и аварийным усечением битых хвостов (\texttt{RecoveryManager}).
    \item Разработан трёхуровневый пользовательский интерфейс, включающий веб-панель оператора на базе \texttt{FastAPI}, \texttt{WebSocket} и \texttt{MJPEG} с интерактивной калибровкой углов доски, выбором камер V4L2/RTSP/файлов, записью и фоновой расшифровкой видеозаписей, пошаговым просмотром партии и неблокирующими окнами разрешения неоднозначности, а также локальный графический экран OpenCV HUD и консольный интерфейс CLI.
    \item Проведено комплексное тестирование и замеры производительности: набор из 285 автоматизированных тестов успешно выполняется с 100\% прохождением; сквозные синтетические бенчмарки подтвердили 100\% точность фиксации стандартных и специальных ходов ФИДЕ (рокировка, взятие на проходе, превращение пешки) при нулевой доле нелегальных ходов и задержке конвейера в пределах бюджета реального времени ($\le 35$~мс/кадр); в среде контейнера с ONNX Runtime измерено время событийного полнокадрового инференса \texttt{ChessQ-Lite INT8} ($910{,}98$~мс), оценки изменений \texttt{DINOv2-Small INT8} ($145{,}77$~мс) и фонового анализа движком Stockfish~16 на глубине 12 ($115{,}19$~мс).
\end{enumerate}

Практический эффект применения разработанной системы заключается в автоматизации цифровой фиксации партий на обычных физических шахматных досках без специализированных электронных сенсоров, снижении трудоёмкости ведения протоколов и обеспечении прозрачного разбора ошибок игроков с объяснимой декомпозицией уверенности каждого зафиксированного хода. К перспективным направлениям дальнейшего развития относятся перенос квантизованных моделей ONNX на встраиваемые граничные устройства класса Jetson/Raspberry~Pi, расширение набора данных реальных нестандартных комплектов фигур для дообучения моделей, а также интеграция мобильного клиента-компаньона.

Личный вклад автора состоит в проведении анализа предметной области и сравнительного обзора существующих подходов, проектировании модульной нейросимвольной архитектуры и контрактов данных, реализации конвейера компьютерного зрения, конечного автомата, логико-символьного скорера, механизма транзакционного журналирования и веб-интерфейса оператора, проведении полного аудита кодовой базы, устранении выявленных дефектов и выполнении серии тестовых и стендовых испытаний.
